\section{Methods II: the modelling framework}
\label{sec:methods-models}

\subsection{Overview}

The framework answers one question: \emph{given how long a person has been unloaded and which muscle is asked
about, how much of that muscle is gone, and how much of the answer depends on anything else?} Two consequences
follow from the wording. The prediction concerns a group, not a person, because the literature reports group
means; and the framework is judged on whether it estimates the duration--muscle relationship honestly, not on how
tightly a flexible model can be made to fit. \Cref{fig:framework} shows the chain from the search to the three
tiers, and \cref{tab:tiers} their roles.

\begin{figure}[htbp]
\centering
\begin{tikzpicture}[node distance=5mm and 5mm]
  \node[box=4.4cm] (search) {\textbf{Literature search}\\ 5,731 records, four sources};
  \node[box=4.4cm, right=of search] (extract) {\textbf{Screening and extraction}\\ frozen schema, verification};
  \node[box=4.4cm, right=of extract] (dataset) {\textbf{Frozen dataset v1.1}\\ 742 rows, 36 campaigns};
  \node[box=4.4cm, above=4mm of extract] (corpus) {Pre-search corpus (9 studies)\\ NASA open data (1 campaign)};
  \node[box=14.4cm, below=7mm of extract] (subset) {\textbf{Modelling subset}: during unloading, lower limb, one tissue per occasion\\
     subset A: 346 rows, 32 campaigns \quad$\cdot$\quad subset B: 304 rows, 25 campaigns \quad$\cdot$\quad exposure: day of the scan};
  \node[keybox=14.4cm, below=5mm of subset] (loco) {\textbf{Leave-one-campaign-out}: 32 folds; no campaign on both sides of a fold, asserted in code;
     intervals by resampling campaigns};
  \node[keybox=4.4cm, minimum height=1.35cm, below=7mm of loco.south west, anchor=north west] (t1) {\textbf{Tier 1}\\ three-level meta-regression\\ \scriptsize coefficients with cluster-robust CIs};
  \node[keybox=4.4cm, minimum height=1.35cm, below=7mm of loco.south, anchor=north] (t2) {\textbf{Tier 2}\\ four ML families vs.\ duration curve\\ \scriptsize + TabPFN post hoc; out-of-campaign error, SHAP};
  \node[keybox=4.4cm, minimum height=1.35cm, below=7mm of loco.south east, anchor=north east] (t3) {\textbf{Tier 3}\\ Jev probabilistic forecast\\ \scriptsize post hoc; ablations, validation};
  \draw[flow] (search) -- (extract);
  \draw[flow] (extract) -- (dataset);
  \draw[flow] (corpus.east) -| (dataset.north);
  \draw[flow] (dataset.south) -- (dataset.south |- subset.north);
  \draw[flow] (subset) -- (loco);
  \draw[flow] (loco.south -| t1.north) -- (t1.north);
  \draw[flow] (loco.south) -- (t2.north);
  \draw[flow] (loco.south -| t3.north) -- (t3.north);
\end{tikzpicture}
\caption[The framework]{The framework. The dataset is frozen before any model is fitted; all three tiers read the
same modelling subset, and tiers 2 and 3 are scored out of campaign against the same duration curve. Tier 1 uses
cluster-robust inference on the campaign rather than cross-validation, because it estimates rather than predicts.}
\label{fig:framework}
\end{figure}

\begin{table}[htbp]
\centering
\caption{The three tiers.}
\label{tab:tiers}
\small
\begin{tabularx}{\linewidth}{@{}l Y Y Y@{}}
\toprule
& Tier 1 --- estimation & Tier 2 --- comparison & Tier 3 --- probabilistic forecast \\
\midrule
Question & What is the duration--response, and which muscles deviate from it? & Does a flexible learner find
  structure the curve misses? & Can a model that reads each campaign's description, and the other campaigns' data,
  do better? \\
Method & Three-level meta-regression, nonlinear in the scan day & Ridge, random forest, SVR, gradient boosting;
  TabPFN added post hoc &
  Language model returning a probability for each range of outcome \\
Validation & Cluster-robust inference on the campaign & Leave-one-campaign-out, nested tuning & Leave-one-campaign-out,
  ablations, recognition probe, validation battery \\
Reports & Coefficients, curve and ranking with 95\% CIs & Out-of-campaign MAE, RMSE, pooled $R^2$ & MAE, CRPS, log
  score, coverage of the most probable ranges \\
Status & Primary result; declared in advance & Declared in advance; allowed to fail & Added after the tier-2 null
  result; not pre-registered \\
\bottomrule
\end{tabularx}
\end{table}

\subsection{Features}
\label{sec:features}

The covariate budget is set by the number of campaigns, not the number of rows. At roughly ten independent studies
per study-level covariate~\citep{higgins2011}, 32 campaigns permit about three study-level terms. Variables that
vary \emph{within} a campaign---muscle, arm, scan day---are far cheaper, because each campaign contributes its own
internal contrast. The implemented feature set is:

\begin{itemize}
  \item \textbf{Days of unloading at the scan}, entered through a nonlinear basis (\cref{sec:duration-forms});
  \item \textbf{Muscle family}, categorical: nine named families plus \code{whole\_limb} in subset A;
  \item \textbf{Arm type}, binary (control or countermeasure). A dose ordinal or a countermeasure category would be
        fitting noise at this sample size: the 126 countermeasure rows are spread across nine named modalities,
        four of them with fewer than ten rows;
  \item \textbf{Modality and outcome}, collapsed to six levels (MRI volume, MRI CSA, CT CSA, DXA lean mass,
        ultrasound CSA, ultrasound thickness), because DXA lean mass and MRI volume do not measure the same thing;
  \item \textbf{Composite}, binary, so that the mixed granularity left by the resolution rule is visible to the model.
\end{itemize}

Age, sex, population, countermeasure modality, measurement site, tilt angle, nutrition control, body mass and the
quality markers are recorded but not modelled: each is sparse, nearly constant, or almost perfectly confounded with
the campaign (for example, 347 of the 425 rows come from men-only groups). They are reserved for stratified
sensitivity analyses. Provenance columns and \code{cohort\_id} are never modelled; the campaign is the grouping
variable, and the code asserts that it never appears as an input. The design document also lists the
percentage-change estimator (\cref{sec:schema}) as a covariate; it is not in the implemented design matrix, and its
effect is to be examined by sensitivity analysis S7 instead. \draftnote{decide whether to add the estimator flag to
the model or to amend the design document.}

\subsection{Tier 1: the three-level meta-regression}
\label{sec:tier1}

\paragraph{Model} For row $i$ measuring muscle $m$ in study $s$ of campaign $c$,
\begin{equation}
  y_{i} = f(t_{i}) + \mathbf{x}_{i}^{\top}\bm{\beta} + u_{c} + v_{s(c)} + w_{m(c)} + \varepsilon_{i},
  \qquad
  \varepsilon_{i} \sim \mathcal{N}\!\left(0, \sigma^2_{\varepsilon}/\omega_{i}\right),
  \label{eq:tier1}
\end{equation}
where $y_i$ is the percentage change, $t_i$ the day of the scan, $f$ one of the duration forms below,
$\mathbf{x}_i$ the categorical and binary features, and $u_c \sim \mathcal{N}(0,\sigma^2_c)$,
$v_{s(c)} \sim \mathcal{N}(0,\sigma^2_s)$ and $w_{m(c)} \sim \mathcal{N}(0,\sigma^2_m)$ random intercepts for the
campaign, the study within the campaign and the muscle within the campaign. The last term is what makes the mixed
granularity legitimate: repeated measurements of the same muscle in the same campaign are correlated, and the model
is told so. The weights $\omega_i$ are the number of participants analysed, normalised to mean one, with missing
values set to the median. This is the standard structure for pooled aggregate data with several effect sizes per
study~\citep{assink2016,harrer2021}, and it does not need the within-study covariances between muscles.

\paragraph{Estimation} Because every random effect is nested inside a campaign, the marginal covariance is block
diagonal, with one block per campaign:
\begin{equation}
  \mathbf{V}_c = \sigma^2_{\varepsilon}\,\mathbf{W}_c^{-1} + \sigma^2_c\,\mathbf{J}_c + \sigma^2_s\,\mathbf{S}_c + \sigma^2_m\,\mathbf{M}_c ,
\end{equation}
where $\mathbf{W}_c$ is the diagonal matrix of weights, $\mathbf{J}_c$ a matrix of ones, and $\mathbf{S}_c$ and
$\mathbf{M}_c$ indicate pairs of rows from the same study and of the same muscle. The fixed effects are profiled out
by generalised least squares at each evaluation, and the four variances are found by maximising the marginal
likelihood on the log scale (L-BFGS-B). Maximum likelihood rather than REML is used because the duration forms do
not share a design matrix and restricted likelihoods cannot be compared across them; the price is a mild downward
bias in the variance components at 32 campaigns. The model was implemented directly in NumPy and SciPy rather than
through a mixed-model library, which carries two levels natively but not the third or the sandwich below.

\paragraph{Inference} Confidence intervals and tests use a cluster-robust sandwich estimator clustered on the
campaign~\citep{liang1986,cameron2015}, on top of the random-effects structure:
\begin{equation}
  \widehat{\operatorname{Cov}}(\hat{\bm\beta}) = \kappa\, \mathbf{A}^{-1}
  \Big( \textstyle\sum_{c} \mathbf{X}_c^{\top}\mathbf{V}_c^{-1}\mathbf{r}_c\,\mathbf{r}_c^{\top}\mathbf{V}_c^{-1}\mathbf{X}_c \Big)
  \mathbf{A}^{-1},
  \qquad
  \mathbf{A} = \textstyle\sum_{c}\mathbf{X}_c^{\top}\mathbf{V}_c^{-1}\mathbf{X}_c ,
\end{equation}
with the small-sample factor $\kappa = \tfrac{G}{G-1}\cdot\tfrac{N-1}{N-p}$ for $G$ campaigns, $N$ rows and $p$
fixed effects. Intervals use a $t$ distribution with $G-1$ degrees of freedom---31 in subset A---never the number
of rows. The random effects model the correlation; the sandwich protects the inference if that model is
misspecified, which at 32 clusters it may be.

\paragraph{Duration forms and the headline rule}
\label{sec:duration-forms}
Three forms of $f$ are fitted and reported side by side:
\begin{itemize}
  \item \textbf{logarithmic}, $f(t) = b\ln t$, the form used in a published synthesis~\citep{marusic2021}, so that the
        two curves can be compared;
  \item \textbf{saturating exponential}, $f(t) = A(1 - e^{-t/\tau})$, whose two parameters are what a
        countermeasure planner wants: how fast the muscle gives up ($\tau$) and how much it will eventually lose
        ($A$, read at the reference scenario). $\tau$ is profiled over the declared grid
        $\{7, 10, 14, 21, 28, 35, 45, 60, 90\}$ days by maximum likelihood and counted as one extra parameter;
  \item \textbf{restricted cubic spline} with three knots at the 10th, 50th and 90th percentiles of the scan day
        (10, 28 and 89 days)~\citep{harrell2015,crippa2019}, as a shape check.
\end{itemize}
The rule, declared before fitting, was: the saturating form is the headline unless another parametric form beats it
by more than 4 AIC points; the spline is never the headline, and serves to show whether the parametric forms miss a
shape in the data.

\paragraph{Reference levels and contrasts} Every contrast is read against a reference level and inherits its
uncertainty, so the reference muscle family was chosen rather than inherited: knee extensors, with 86 rows across 22
campaigns in subset B, against the alphabetical default of dorsiflexors with 8 rows from 4 campaigns. The
antigravity comparison---plantar flexors against dorsiflexors---is computed as a separate contrast from the
coefficient difference, so that the choice of reference cannot change it. The modality reference is the
alphabetical first level, CT cross-sectional area. The curve is drawn for one scenario: a control arm, the
reference muscle family and modality, not a composite. Other scenarios shift the curve by their coefficients; they
do not reshape it.

\paragraph{Fits} The model is fitted twice with the same specification: on subset A for the duration--response and
on subset B for the muscle ranking, which is reported as the fitted change at day 60 for a control arm.

\subsection{The duration-only baseline}
\label{sec:baseline}

The baseline is the simplest honest answer to ``how much muscle is gone after $t$ days'': percentage change as a
function of the scan day and nothing else, in linear, logarithmic and saturating form, fitted by weighted least
squares (weights $\omega_i$) inside the same leave-one-campaign-out loop as everything else, with $\tau$ chosen from
the same grid within each training fold. The best-scoring form is the reference against which tier 2 is judged.

\subsection{Tier 2: the comparative machine-learning families}
\label{sec:tier2}

Four families were fitted on subset A, as promised in the accepted abstract: ridge regression (penalised linear
regression)~\citep{hoerl1970}, random forest~\citep{breiman2001}, support vector regression with a radial-basis
kernel~\citep{smola2004} and gradient boosting~\citep{friedman2001}, all through scikit-learn~\citep{pedregosa2011}.
The input is the feature set of \cref{sec:features} with the duration entered as the saturating basis at a fixed
$\tau$ of 21 days; the flexible families can reshape it.

\paragraph{Validation} The fold unit is the campaign: 32 folds, each training on 31 campaigns and predicting one it
has never seen. A random or paper-level split would place the same participants on both sides of the fold
boundary and measure memorisation. The leakage guard is an assertion, not a convention: for every fold the code
asserts that the training and test campaign sets are disjoint and that the campaign identifier is absent from the
design matrix, and a test deliberately feeds it a random split to confirm that it fails.

\paragraph{Nested tuning} Hyperparameters were chosen inside each training fold by a grid search with five inner
folds grouped by campaign, scored by mean absolute error, so that the outer fold sees none of the choice. The grids
were small on purpose: ridge penalty $\alpha \in \{0.1, 1, 10, 100\}$ on standardised inputs; random forest (500
trees) minimum leaf size $\in \{2, 4, 8\}$ and feature fraction $\in \{0.5, 1\}$; SVR $C \in \{1, 10, 100\}$ and
$\gamma \in \{0.05, 0.1, 0.5\}$ on standardised inputs; gradient boosting (300 trees) depth $\in \{1, 2\}$ and
learning rate $\in \{0.03, 0.1\}$. The tier-2 families were fitted unweighted.

\paragraph{Metrics} Mean absolute error (MAE) and root mean squared error are computed within each held-out
campaign and averaged with one vote per campaign, so that the largest campaign cannot dominate. The 95\% interval is
a percentile bootstrap over the 32 campaign-level errors (2,000 replicates, seed 20260918). $R^2$ is computed once
over all pooled out-of-fold predictions rather than averaged over folds: most campaigns ran a single duration, so a
duration-only model predicts one value across a held-out fold and scores below that fold's own mean by construction,
whatever the quality of the curve. The first baseline run returned about $-8.5$ per fold and $+0.06$ pooled from
identical predictions.

\paragraph{A fifth family, post hoc: TabPFN} On 19 September, after the null result of the four families was
known, a fifth was added: TabPFN~\citep{hollmann2025}, a transformer pretrained on millions of synthetic tabular
datasets that predicts in context, conditioning on the training rows it is given rather than fitting weights to
them. It is the neural network built for exactly this regime---tables of a few hundred rows---and, having no
hyperparameters to tune, it cannot be tuned against the outer fold. It was run with fixed declared defaults (version
2.0.9, on CPU, seed 20260918) on the same design matrix, the same 32 folds and leakage assertions and the same
campaign-weighted scoring, against the same logarithmic duration curve. Because tabpfn 2.0.9 requires scikit-learn
below 1.7 while the four families were fitted with 1.8, it runs in a separate environment (scikit-learn 1.6.1,
PyTorch 2.14) through its own runner, \file{framework/run\_tabpfn.py}, and writes its own results file; the four
families' results are untouched. Like tier~3, it is reported as post hoc. An earlier run of the same family, on
\code{dataset\_v1.0} and the campaigns' planned lengths, is superseded.

\paragraph{Importance} SHAP values~\citep{lundberg2017} were computed exactly (tree explainer) for the best-scoring
family, recomputed inside every fold, and summarised as the share of folds in which each feature reaches the top
three. Stability is reported first and magnitude second; nothing about importance is read causally.

\subsection{Tier 3: a probabilistic language-model forecast}
\label{sec:tier3}

\begin{keybox}{Status of tier 3}
Tier 3 was added on 21 September 2026, after the tier-2 null result was known. It is not pre-registered. What could
still be done was to declare the method and every bar before the first answer was requested, and to commit that
declaration to the repository first; the same was done for the ablations and for the validation battery, each
before its own first answer.
\end{keybox}

The question tier 2 leaves open is whether something that reads the \emph{description} of a campaign---who the
participants were, what they did, which muscle, which scanner---can extract what the coded columns of tier 2 could not.
Tier 3 uses TypeSafe's \Jev{} (\code{jev-1.13.0}, pinned)~\citep{typesafe2026}, a model built to return a
calibrated probability for each option of a fixed set given a described state. It is never asked for a number. For
each row it is asked one question---\emph{which of these ranges will the measurement fall in?}---and every numeric
quantity is computed in code from the probabilities it returns.

\paragraph{Two arms}
\begin{itemize}
  \item \textbf{Without history} (346 rows, 32 campaigns): the model sees nothing of the held-out campaign, as the
        tier-2 families do, and forecasts the percentage change from baseline over 2\pp{} ranges from $-30$ to
        $+4$, open at both ends. Reference: the logarithmic duration curve fitted to the training campaigns.
  \item \textbf{With history} (84 rows, 5 campaigns): the model also sees the held-out campaign's scans strictly
        before the target day and forecasts the change since the previous scan, over 1\pp{} ranges from $-14$ to
        $+4$. References: no further change (\emph{last scan}) and the curve's step from the last scan's day to the
        target day (\emph{last scan plus curve step}). One campaign, BBR1, carries 39 of the 84 rows.
\end{itemize}

\paragraph{What the model is shown} A short background paragraph; the participants, protocol, countermeasure and
measurement of the target row in words; the target day and its distance from the planned end; the duration curve
fitted to the training campaigns; and as many of the training campaigns' observations as fit a budget of 26,000
tokens, most relevant first (same muscle, then same family, then same kind of group, then nearest day). The budget
was set from a measured rate of 2.3 characters per token so that state and question fit the model's 32,000-token
limit.

\paragraph{What it is never shown, by assertion} Any row of the held-out campaign in the reference tables; any
scan of the held-out campaign on or after the target day; the target's own outcome; and anything that names a
paper, an author, a campaign or a registry. The last rule also keeps the model from recognising a published study by
name and recalling its number. The assertions run on every request.

\paragraph{Baselines as distributions} Each baseline's forecast distribution is its point prediction plus the
residuals it made on the training campaigns, binned into the same ranges. The model and the baselines are therefore
scored on identical footing for every metric.

\paragraph{Scores} The point forecast is the probability-weighted mean of the ranges' representative values (the
midpoint, or half a step beyond an open end). The continuous ranked probability score (CRPS) is computed as the
ranked probability score~\citep{epstein1969} multiplied by the range width, so that it reads in percentage points
and reduces to the absolute error for a forecast with no spread~\citep{gneiting2007}. The log score is minus the
log of the probability given to the range holding the truth, floored at $10^{-6}$. The \emph{most probable range}
at level $\alpha$ is the shortest run of adjacent ranges whose probability reaches $\alpha$; its coverage and width
are reported at 50\% and 80\%. Every metric is averaged within a held-out campaign and then across campaigns. The
comparison statistic is the \emph{paired gain}: the reference's campaign MAE minus the model's, averaged over
campaigns, with a campaign bootstrap interval.

\paragraph{Bars, declared before the first answer}
\begin{itemize}
  \item MAE: at least 15\% relative improvement over the arm's reference; otherwise report a null result.
  \item CRPS: at least 15\% relative improvement, reported alongside the MAE and never instead of it.
  \item Coverage of the 80\% range between 70\% and 90\%; otherwise the ranges may not be called 80\% ranges.
  \item Leakage assertions hold for every request, and \code{make forecast} rebuilds every number from the committed
        answer cache without calling the service. Both blocking.
\end{itemize}

\paragraph{What tier 3 cannot rule out} \Jev{} was pretrained on text that may include some of these papers.
Withholding every name is the strongest guard available, but a model that has read a paper could still recognise it
from a description plus a set of numbers. The ablations below were designed for that concern.

\subsubsection{Ablations and the recognition probe}
\label{sec:ablation-methods}

Three ablations on the arm without history each change one thing about what the model is shown, and are scored
exactly as the first run was, against the same duration curve:
\emph{generic}, in which the target is described only by muscle, role, imaging method, kind of group and day---no
participants, protocol text, planned length or measurement site, roughly the information tier 2 had;
\emph{scrambled reference}, in which the other campaigns' values are shuffled among their rows; and
\emph{no reference}, with no curve and no other campaigns' observations.

The \emph{recognition probe} shows the model the target's full description, exactly as the forecast sends it, and
asks which of the 11 campaigns in subset A that have a proper name it comes from (AGBRESA, BBR1, BBR2-2, BRACE,
LunHab, MEDES LTBR, NASA SPRINT, UTMB Campaign 3, PlanHab, VBR, WISE-2005), or none of them. Chance is one in
twelve.

The reading rules were declared before any answer: a campaign is \emph{recognised} if the model's top choice is its
true name on at least half of its rows; recognition is a \emph{live explanation} of the gain if two or more of the
three campaigns carrying most of the gain (NASA SPRINT, BBR2-2, MEDES LTBR) are recognised, or if the Spearman
correlation between a campaign's mean probability on its true name and its gain is 0.5 or more; the generic variant
\emph{keeps the gain} if its paired gain has an interval above zero and is at least half the full run's (0.21\pp);
and the scrambled variant \emph{shows that the model reads the reference data} if its paired-gain interval reaches or
falls below zero. If the generic variant keeps the gain and recognition is not live, the point forecasts may be
quoted as a result.

\subsubsection{The validation battery}
\label{sec:validation-methods}

A second set of runs, also declared before its first answer and also leave-one-campaign-out, asks whether the result
depends on how the question was put and whether the history arm reads what it is given. On both arms,
\emph{reordered} shows the same reference rows in a different order and \emph{shifted ranges} moves every range edge
by half a step; neither changes what the model knows. On the history arm, the three ablations are repeated and a
fourth, \emph{scrambled history}, shuffles the held-out campaign's earlier values among its earlier scans. A
\emph{repeat check} sends the first 20 requests of the arm without history again, past the cache. A result is
\emph{robust to presentation} if under both presentation checks its MAE moves by less than 10\% and its paired gain
keeps its sign (with an interval above zero on the arm without history); the model is \emph{consistent} if all 20
repeats return identical probabilities; and the model \emph{reads the earlier scans} if scrambling them worsens the
error with an interval below zero. Tier 3 is called working rather than lucky if the arm without history is robust
and consistent and the history arm reads what it is given.

\subsection{Success criteria declared before modelling}
\label{sec:criteria}

The thresholds in \cref{tab:criteria} were agreed in the project plan before any model was fitted. The first three
were deliberately allowed to fail; the last three were blocking. They are assessed in \cref{sec:results-criteria}.

\begin{table}[htbp]
\centering
\caption{Success criteria for tiers 1 and 2, declared before modelling (project plan, section 8).}
\label{tab:criteria}
\small
\begin{tabularx}{\linewidth}{@{}L{4.6cm} L{3.6cm} Y@{}}
\toprule
Criterion & Target & If missed \\
\midrule
Out-of-campaign MAE & $\leq$ 4\pp & Report honestly; the dataset and framework remain the contribution \\
Best model vs.\ duration-only baseline & $\geq$ 15\% relative MAE improvement & Report the null result \\
SHAP top-3 stability & Same top three in $\geq$ 80\% of folds & Present importance as indicative only \\
Leakage audit & No campaign split across folds; no preprocessing outside a fold & Blocking \\
Reproducibility & One command regenerates every number and figure & Blocking \\
Physiological sign-off & Scientific lead signs off the fitted relationships & Blocking \\
\bottomrule
\end{tabularx}
\end{table}

\subsection{Sensitivity analyses}
\label{sec:sensitivity-methods}

Eleven sensitivity analyses were declared, each the same pipeline with one change (\cref{tab:sensitivity}). S1, S2
and S4 were pre-registered as analyses that must be shown whatever they say. At the time of writing only S6 has been
run (\cref{sec:results-s6}).

\begin{table}[htbp]
\centering
\caption{Declared sensitivity analyses and their status on 21 September 2026. $\star$ = must be shown in the report.}
\label{tab:sensitivity}
\small
\begin{tabularx}{\linewidth}{@{}l Y Y l@{}}
\toprule
& Change & What it tests & Status \\
\midrule
S1$^\star$ & Composite-first instead of component-first resolution & Whether the ranking is an artefact of granularity & Not run \\
S2$^\star$ & Drop the MEDES 90-day campaign & Whether one campaign carries the duration--response & Not run \\
S3 & Drop the two Berlin campaigns & The same for the other dominant group & Not run \\
S4$^\star$ & MRI volume only & Whether mixing modalities changed the answer & Not run \\
S5 & Exclude figure-derived and low-confidence rows & Whether digitisation is load-bearing & Not run \\
S6 & Inverse-variance weights & Whether the weighting scheme matters & \textbf{Run} \\
S7 & Each percentage estimator alone & Whether the two estimators agree & Not run \\
S8 & Add mean age as a study-level covariate & Whether the older cohorts shift anything & Not run \\
S9 & Sex-stratified, if the folds survive it & A question that will be asked & Not run \\
S10 & Recovery rows with a phase term & Whether reconditioning can be modelled & Not run \\
S11 & Rows with a stated measurement site, then by position & Whether slice position moves the estimate & Not run \\
\bottomrule
\end{tabularx}
\end{table}

\subsection{Software and reproducibility}
\label{sec:reproducibility}

The framework is a set of single-purpose Python modules (\file{framework/}) driven by one configuration file,
\file{framework/config.yaml}, which holds every subset predicate, feature list, duration form, model grid,
bootstrap setting and the random seed; no module contains a number that is not in the configuration. The core---
loading, features, folds, baseline, evaluation and tier 1---needs only NumPy, pandas, SciPy and PyYAML, so that the
primary result does not depend on the optional stack that produced the null result. The results reported here were
produced with Python 3.13.0, NumPy 2.4.6, pandas 2.2.3, SciPy 1.17.0 and scikit-learn 1.8.0, and each results file
records the dataset version, its hash and the commit it came from.

\code{make all} regenerates every result from the frozen dataset. Tier 3 is rebuilt from the model's answers
committed in \file{results/forecast\_cache/} (430 answers for the first run, 1,280 for the ablations and 1,195 for
the validation battery), never by calling the service again. The test suite holds 234 checks in 20 files, including
the deliberately failing leakage test. The post-hoc TabPFN family is rebuilt by \code{make tabpfn} in its own
environment and is not part of \code{make all}. The figures and tables of this report are regenerated from the
results files by \file{report/make\_figures.py}.
